\section{Reasoning: de "saber pensar" a "saber hacer"}

\subsection{El objetivo de ingeniería de Reasoning}
Yu Su insistió repetidamente en el curso: el objetivo final de Reasoning no es exhibir la cadena de razonamiento, sino mejorar la calidad de la acción (Action Quality). Si el razonamiento no produce mejores secuencias de acciones, una tasa de error más baja o una mayor eficiencia de muestreo, entonces ese razonamiento es un gasto inútil.

\begin{figure}[H]
\centering
\includegraphics[width=0.84\textwidth]{slide-02-reasoning-action.jpg}
\caption{Diapositiva hacia el minuto 00:21:50 del video: el acoplamiento entre Reasoning y Acting es el eje del curso}
\end{figure}

\begin{importantbox}{Reasoning for Acting}
\begin{itemize}
    \item Usar el razonamiento para mejorar la calidad de la selección de acciones, no para aumentar la extensión superficial del pensamiento;
    \item El razonamiento debe poder desencadenar la corrección de la estrategia (replan) y el cambio de herramienta (tool switch);
    \item El valor del razonamiento se refleja, en última instancia, en tres indicadores: tasa de éxito, costo y robustez.
\end{itemize}
\end{importantbox}

\subsection{De CoT al razonamiento ejecutable}
El curso coloca CoT, ReAct, Self-Reflection y Meta-Reasoning en un mismo espectro continuo: en esencia, todos consisten en "inyectar pensamiento estructurado antes y después de la acción". La diferencia no está en si hay razonamiento o no, sino en cómo el razonamiento conecta el estado del entorno con la interfaz de acción.

\begin{knowledgebox}{Cuatro tipos comunes de acciones de razonamiento}
\begin{enumerate}
    \item Descomposición de tareas: dividir un objetivo complejo en subobjetivos verificables;
    \item Alineación de evidencia: corregir las hipótesis internas con salidas externas de recuperación o herramientas;
    \item Atribución de fallas: identificar el origen del error (falta de conocimiento, falla de herramienta, error de planeación);
    \item Reescritura de estrategia: reescribir el plan de acciones posteriores con base en el resultado de la atribución.
\end{enumerate}
\end{knowledgebox}

\begin{warningbox}{Más tokens de razonamiento no siempre es mejor}
Sin un verificador eficaz ni restricciones de límite, una cadena de razonamiento larga puede provocar:
\begin{itemize}
    \item Explosión de costos;
    \item Amplificación de la alucinación de confianza;
    \item Que las rutas erróneas se "expliquen de un modo que parezca más correcto".
\end{itemize}
\end{warningbox}

\subsection{Resumen de la sección}
Reasoning es el mecanismo de ganancia en la decisión del Agente. Lo que realmente importa no es "si hay razonamiento o no", sino "si el razonamiento cambió la acción hacia una mejor".
